\chapter{Transactions de portefeuille et risk bids}\label{ch:mx:portfolio-trades-and-risk-bids}

Un fonds de pension demande à trois courtiers un prix auquel ils accepteront de prendre en charge un panier de cinq cents actions à la clôture de ce soir, en ne connaissant que sa taille,
ses secteurs et sa liquidité. Le courtier qui l'emporte apprend ensuite les noms et découvre s'il a offert trop cher. Ce chapitre évalue une telle offre sur
un panier simulé~: le coût de sa liquidation, le risque porté pendant cette liquidation, et la minoration que la concurrence impose à un enchérisseur qui ne veut pas
gagner seulement lorsqu'il se trompe~; il compare ensuite un panier divulgué à un panier à l'aveugle et se termine par une transition entre deux portefeuilles.

\section{Transactions de portefeuille en agence et en principal}

\begin{definition}[Transaction de portefeuille en agence, risk bid]\label{def:mx:portfolio-trades-and-risk-bids:agency}
Une \emph{transaction de portefeuille en agence}\index{transaction de portefeuille en agence} exécute une liste d'ordres pour un client en tant qu'agent, le client supportant les coûts et le
risque et versant une commission. Un \emph{risk bid}\index{risk bid} est un prix en principal pour la liste entière~: le courtier achète (ou vend) chaque position à
un indice de référence convenu, en général la clôture, diminué d'une décote, et supporte le coût et le risque de leur débouclage.
\end{definition}

\begin{definition}[Risk bid à l'aveugle]\label{def:mx:portfolio-trades-and-risk-bids:blind}
Un \emph{risk bid à l'aveugle}\index{risk bid à l'aveugle} est un \omterm{def:mx:portfolio-trades-and-risk-bids:agency}{risk bid} formulé à partir d'une description du panier (sa taille, ses secteurs, sa liquidité et ses caractéristiques
de risque), sans les noms, qui ne sont révélés qu'au gagnant une fois la transaction conclue.
\end{definition}

Une transaction de portefeuille (le livre~2, chapitre~22) déplace une liste entière d'un seul coup~: un rééquilibrage d'indice, des entrées de fonds, un changement de gérant. Le client
choisit entre payer lui-même les coûts (en agence) et en confier la charge à un courtier (en principal)~; le \omterm{def:mx:portfolio-trades-and-risk-bids:agency}{risk bid} se valorise comme tout transfert de risque
(chapitre~21), avec deux particularités~: la liste est large et diversifiée, de sorte que son risque peut être couvert en partie, et, dans la version à l'aveugle, l'enchérisseur valorise
une description. Kavajecz et Keim (2005) ont étudié de telles enchères en principal à l'aveugle et constaté qu'elles offraient au demandeur de liquidité une économie de coûts de transaction
par rapport à un trading plus traditionnel, et aux fournisseurs de liquidité un moyen efficace d'obtenir du flux d'ordres.

Le panier simulé provient du marché synthétique du livre~7 (\texttt{firm.synthmkt})~: 786 titres cotés le dernier jour, avec un volume quotidien médian de
37.5~millions de USD et une volatilité quotidienne médiane de 171 points de base~; les demi-fourchettes suivent une règle de taille (une hypothèse~: une fourchette complète médiane de 5.9 points
de base), et un modèle de risque statistique à dix facteurs sur 250 jours de rendements fournit la covariance (\texttt{firm.riskmodel} du livre~7). Le panier du fonds de pension
compte 500 de ces titres, pour 250~millions de USD, avec des positions allant de 38\,000~USD à 4.8~millions de USD, une médiane de 1.0~\% du volume quotidien.

\section{Le coût de liquidation d'un panier}

\begin{definition}[Profil de liquidité du panier]\label{def:mx:portfolio-trades-and-risk-bids:profile}
Un \emph{profil de liquidité du panier}\index{profil de liquidité du panier} résume un panier sans ses noms~: la part de son notionnel dans des classes de taille de position
rapportée au volume quotidien, par secteur, avec son bêta, sa volatilité et ses autres caractéristiques de risque.
\end{definition}

Chaque position est vendue à hauteur de 20~\% au plus de son volume quotidien, de sorte qu'une position de $k$ jours de volume prend $\max(1,5k)$ jours~; son coût est sa
demi-fourchette, ajoutée au terme $Y\sigma\sqrt{\min(Q/V,0.2)}$ de la loi en racine carrée (chapitre~11, $Y=0.7$). Le \omterm{def:mx:portfolio-trades-and-risk-bids:profile}{profil de liquidité du panier} figure dans la
\cref{fig:mx:portfolio-trades-and-risk-bids:profile}~: 98.8~\% de son notionnel est inférieur à 5~\% du volume d'une journée, et aucune position n'exige plus d'un jour. Le
coût de débouclage du panier pour le gagnant s'élève à 14.9 points de base du panier (\texttt{firm\_riskbid.liquidation}).

\begin{omfigure}
\begin{tikzpicture}
\begin{axis}[width=9.4cm,height=5.2cm,ybar,bar width=14pt,ylabel={\small part du notionnel du panier (\%)},
  xtick={0,1,2,3,4,5},xticklabels={{$<0.5~\%$},{0.5--1~\%},{1--2~\%},{2--5~\%},{5--10~\%},{10--20~\%}},
  xlabel={\small position rapportée au volume quotidien},tick label style={font=\footnotesize},ymin=0,ymax=40,
  grid=major,grid style={gray!25},enlarge x limits=0.12,table/col sep=comma]
\addplot[fill=omDef!70,draw=omDef] table[x=i,y=share]{figdata/microstructure/22-portfolio-trades-and-risk-bids/profile.csv};
\end{axis}
\end{tikzpicture}

\omcaption{Le \omterm{def:mx:portfolio-trades-and-risk-bids:profile}{profil de liquidité du panier}~: son notionnel par taille de position rapportée au volume quotidien, ce que voit un enchérisseur à l'aveugle (avec les poids
sectoriels). Données~: \texttt{mx\_riskbid.blind}.}\label{fig:mx:portfolio-trades-and-risk-bids:profile}
\end{omfigure}

\section{Valorisation d'un risk bid}

Pendant qu'il vend, le gagnant conserve ce qui reste. Chaque position étant vendue progressivement, de façon linéaire, sur ses jours $d_i$, la variance du P\&L vaut
$\sum_{ij}w_iw_jC_{ij}I(d_i,d_j)$, où $C$ est la covariance quotidienne et $I(a,b)=a/2-a^2/(6b)$ pour $a\le b$ le recouvrement des deux détentions (pour un seul jour,
un tiers de la variance d'une journée)~: c'est le terme de risque d'Almgren--Chriss multi-actifs (Almgren et Chriss, 2001~; chapitre~14) pour les calendriers linéaires. Pour
le panier, un écart type vaut 40.4 points de base~; vendu à découvert sur contrats à terme sur indice dès le départ pour supprimer son bêta de marché, il vaut 12.5.

\omcode{code/firm/riskbid/firm_riskbid.py}{43}{57}{Le risque de détenir un panier pendant qu'on le vend progressivement~: chaque paire de positions contribue par sa covariance multipliée par la durée de recouvrement de leurs détentions linéaires~; une couverture de marché retire la part du facteur dans la covariance.}\label{lst:mx:portfolio-trades-and-risk-bids:risk}

L'offre, exprimée comme une décote par rapport à la clôture, additionne quatre composantes~: le coût espéré, une charge pour le risque (ici un écart type du P\&L couvert, le
choix du desk), la minoration qu'impose la concurrence (section suivante), et une marge (2 points de base). Pour le panier divulgué~: 14.9 de coût,
12.5 de risque, 4.2 de minoration et 2.0 de marge, soit une offre à 33.6 points de base sous la clôture (\cref{fig:mx:portfolio-trades-and-risk-bids:bids}).

\section{Risk bids à l'aveugle et malédiction du vainqueur}

Chaque enchérisseur estime le même coût avec sa propre erreur~; celui dont l'estimation est la plus basse remporte l'enchère, et il la remporte le plus souvent lorsque son estimation est trop basse.
C'est la malédiction du vainqueur d'une enchère à valeur commune (le livre~2, chapitre~30). Si chacun des quatre enchérisseurs (le desk et trois concurrents) estime
le coût avec une erreur indépendante d'écart type $s$, l'estimation du gagnant est en moyenne trop basse de $s\,\E[\max\text{ de 4 normales centrées réduites}]=1.03s$
et un enchérisseur qui ne minore pas de ce montant perd sur les transactions qu'il remporte (\texttt{firm\_riskbid.curse}).

\omcode{code/firm/riskbid/firm_riskbid.py}{69}{80}{La malédiction du vainqueur et l'enchère qui la vérifie~: la minoration est l'avance espérée du meilleur de n + 1 estimations~; l'enchère simulée renvoie le profit moyen du gagnant.}\label{lst:mx:portfolio-trades-and-risk-bids:curse}

Avec des erreurs d'estimation de 15~\% du coût et du risque (4.1 points de base pour le panier divulgué), la minoration vaut 4.2 points de base~; sur 200\,000
enchères simulées, les enchérisseurs qui n'ajoutent que la marge de 2 points de base réalisent $-2.2$ points de base sur les transactions qu'ils remportent, et ceux qui minorent aussi
obtiennent les 2.0 qu'ils visaient.

Un enchérisseur à l'aveugle ne voit que le profil. Pour le valoriser, le desk tire 300 paniers ayant le même profil (chaque position conserve son secteur, son notionnel et
sa classe de jours de volume, avec un nom tiré parmi ceux qui conviennent) et les valorise un par un~: leur coût moyen vaut 17.0 points de base (écart type 0.4)
et leur risque couvert 13.0 (1.1). Le profil ne révèle pas que le fonds détenait davantage de titres liquides qu'un tirage au hasard ne le donnerait~: l'estimation à l'aveugle
dépasse de 2.1 points de base le coût divulgué. L'erreur d'estimation plus grande ajoute 0.2 à la minoration. L'offre à l'aveugle s'élève à 36.4 points de base, soit 2.8 de plus que
l'offre divulguée~; le client paie la différence pour ne pas montrer ses noms avant la transaction.

\begin{omfigure}
\begin{tikzpicture}
\begin{axis}[width=10.4cm,height=4.2cm,xbar stacked,bar width=13pt,xlabel={\small offre sous la clôture (points de base)},
  ytick={0,1},yticklabels={à l'aveugle,divulguée},tick label style={font=\footnotesize},xmin=0,xmax=40,ymin=-0.6,ymax=1.6,
  legend columns=4,legend style={font=\footnotesize,at={(0.5,-0.42)},anchor=north,draw=none,fill=none,
  /tikz/every even column/.append style={column sep=6pt}},table/col sep=comma]
\addplot[fill=omDef!70,draw=omDef] table[y=y,x=cost]{figdata/microstructure/22-portfolio-trades-and-risk-bids/bids.csv};
\addplot[fill=omThm!60,draw=omThm] table[y=y,x=risk]{figdata/microstructure/22-portfolio-trades-and-risk-bids/bids.csv};
\addplot[fill=omProp!70,draw=omProp] table[y=y,x=curse]{figdata/microstructure/22-portfolio-trades-and-risk-bids/bids.csv};
\addplot[fill=gray!40,draw=gray] table[y=y,x=margin]{figdata/microstructure/22-portfolio-trades-and-risk-bids/bids.csv};
\legend{coût,charge de risque,malédiction du vainqueur,marge}
\end{axis}
\end{tikzpicture}

\omcaption{Les offres divulguée et à l'aveugle sur le même panier, composante par composante (points de base du panier sous la clôture)~: le coût et le risque de l'enchérisseur à l'aveugle
sont ses estimations à partir du profil de liquidité. Données~: \texttt{mx\_riskbid.bids}.}\label{fig:mx:portfolio-trades-and-risk-bids:bids}
\end{omfigure}

L'offre à l'aveugle conserve néanmoins une utilité~: les noms du client ne fuient pas vers les perdants, qui pourraient sinon négocier avant le débouclage du gagnant (la \omterm{def:mx:dark-pools-and-blocks:leakage}{fuite
d'information} du chapitre~9). Savoir si les 2.8 points de base achètent une protection suffisante dépend du nombre de dealers qui verraient les noms et de ce qu'ils en feraient.

\section{Transitions}

Une transition déplace un fonds d'un portefeuille (celui d'un gérant qu'il quitte) vers un autre. Le gestionnaire de transition (livre~8, chapitre~29) compense les noms
que les deux portefeuilles ont en commun, croise en interne ce qui peut l'être contre d'autres flux, et négocie le reste. Pour deux portefeuilles de 250~millions de dollars américains comptant 150 noms
en commun, vendre l'un et acheter l'autre fermes négocie 500~millions de dollars américains et coûte 50.4 points de base du fonds~; compenser les noms communs négocie
322~millions de dollars américains à 32.8~; croiser en interne 20~\% du reliquat laisse 258~millions de dollars américains à négocier, à 24.0 (\texttt{firm\_riskbid.transition}).
Le croisement interne (livre~7, chapitre~27) et le central risk book (livre~9, chapitre~25) relèvent de la même idée, à l'échelle d'une banque~:
la transaction la moins coûteuse est celle qui n'atteint jamais le marché.

\section{Tutoriel~: tarification d'un risk bid}\label{tut:mx:portfolio-trades-and-risk-bids:tut}

\begin{tutorial}
\textbf{Objectif.} Tarifer une offre pour compte propre sur un panier de 500 noms, divulguée et à l'aveugle, avec la décote liée à la malédiction du vainqueur, et chiffrer une transition.
\textbf{État final~:} les \cref{fig:mx:portfolio-trades-and-risk-bids:profile,fig:mx:portfolio-trades-and-risk-bids:bids} et les chiffres des sections 2 à 5.
\begin{tutsteps}
\item \textbf{Univers.} \texttt{mx\_riskbid.universe()} sur \texttt{firm.synthmkt} et \texttt{firm.riskmodel.pca\_model}~; \texttt{basket()}.
\item \textbf{Coût et risque.} \texttt{firm\_riskbid.liquidation}, \texttt{risk\_sd}~; \texttt{price(names, notional)}.
\item \textbf{À l'aveugle.} \texttt{profile}~; \texttt{blind()}.
\item \textbf{Offre.} \texttt{curse}, \texttt{auction}~; \texttt{bids()}~; \texttt{transition\_study()}~; tracer avec \texttt{fig\_riskbid.py}.
\end{tutsteps}
\textbf{Pour aller plus loin.} Ajouter un \omterm{def:mx:metaorders-latent-liquidity-and-cross-impact:cross}{impact croisé} entre les noms (chapitre~13)~; laisser corrélées les erreurs des concurrents (un modèle commun)~; rendre
le panier moins liquide que ne le suggère son profil.
\end{tutorial}

\section{Construction~: risk bids}\label{bld:mx:portfolio-trades-and-risk-bids:riskbid}

\begin{build}
\bfield{Rôle} La tarification des paniers en principal et des transitions pour le desk du chapitre~20 et les desks bancaires du livre~9.
\bfield{Interface} \texttt{liquidation(notional, adv, sigma, spread\_bp, cap, y)}, \texttt{risk\_sd(notional, cov, days, hedge)},
\texttt{profile(notional, adv, sector, edges)}, \texttt{curse(sigma\_est, competitors)},
\texttt{auction(true\_cost, sigma\_est, bidders, shade, margin, trials, seed)}, \texttt{transition(old, new, cross\_share)}.
\bfield{Règles} Coûts en points de base du notionnel brut du panier~; liquidation linéaire sur les jours de chaque nom~; une couverture retire un facteur
d'un seul coup~; erreurs des enchérisseurs indépendantes.
\bfield{Tests d'acceptation} \texttt{code/firm/riskbid/tests/}~: la règle de liquidation~; le profil~; le risque linéaire comparé à une intégrale numérique
et à une couverture parfaite~; la décote de malédiction fait égaler le profit du gagnant à la marge~; la compensation et le croisement de la transition.
\bfield{Pour aller plus loin} \omterm{def:mx:metaorders-latent-liquidity-and-cross-impact:cross}{Impact croisé}~; enchérisseurs corrélés~; calendriers de liquidation optimaux au lieu de trajectoires linéaires.
\end{build}

\begin{omsources}
\item R.~Almgren and N.~Chriss, ``Optimal execution of portfolio transactions'', \emph{Journal of Risk} 3(2), 2001.
\item K.~A.~Kavajecz and D.~B.~Keim, ``Packaging liquidity: blind auctions and transaction efficiencies'', \emph{Journal of Financial and
Quantitative Analysis} 40(3), 2005.
\end{omsources}

\section{Exercices}

\begin{exercise}[$\star$]\label{exo:mx:portfolio-trades-and-risk-bids:1}
Une position représentant 30~\% du volume d'une journée sur une action dont la volatilité quotidienne est de 200 points de base et la fourchette de 6 points de base~: combien de jours faut-il pour la liquider à
un cap de 20~\%, et quel en est le coût~?
\end{exercise}

\begin{exercise}[$\star$]\label{exo:mx:portfolio-trades-and-risk-bids:2}
Une position vendue de façon linéaire sur trois jours, avec une volatilité quotidienne de 150 points de base~: quel est l'écart type de son P\&L, en
points de base de la position~?
\end{exercise}

\begin{exercise}[$\star$]\label{exo:mx:portfolio-trades-and-risk-bids:3}
Quelle décote un enchérisseur doit-il appliquer face à cinq concurrents dont l'erreur d'estimation est de 4 points de base~?
\end{exercise}

\begin{exercise}[$\star\star$]\label{exo:mx:portfolio-trades-and-risk-bids:4}
Pourquoi la couverture à l'aide de contrats à terme sur indice ramène-t-elle le risque du panier de 40.4 à 12.5 points de base, et pourquoi pas davantage~?
\end{exercise}

\begin{exercise}[$\star\star$]\label{exo:mx:portfolio-trades-and-risk-bids:5}
Pourquoi l'enchérisseur à l'aveugle surestime-t-il le coût de ce panier, et pourrait-il sous-estimer celui d'un autre~?
\end{exercise}

\begin{exercise}[$\star\star$]\label{exo:mx:portfolio-trades-and-risk-bids:6}
Si les quatre enchérisseurs utilisent le modèle de coût d'un même fournisseur, que devient la malédiction du vainqueur~?
\end{exercise}

\begin{exercise}[$\star\star\star$]\label{exo:mx:portfolio-trades-and-risk-bids:7}
\emph{Programmation.} Tarifer le panier divulgué avec une charge de risque de deux écarts types au lieu d'un. Quelle est l'offre~?
\end{exercise}

\begin{exercise}[$\star\star\star$]\label{exo:mx:portfolio-trades-and-risk-bids:8}
\emph{Trouver la faille.} «~Nous avons remporté 40~\% des \omterm{def:mx:portfolio-trades-and-risk-bids:agency}{risk bids} que nous avons faits l'an dernier, donc notre tarification est à peu près juste.~»
\end{exercise}

\section{Problème~: tarifer le panier à l'aveugle}

\begin{problem}[{Problème du week-end --- tarifer le panier à l'aveugle}]\label{pb:mx:portfolio-trades-and-risk-bids:1}
Un fonds de pension propose un panier de 500 noms pour une offre pour compte propre à la clôture~; le desk est en concurrence avec trois autres.

\medskip\noindent\textbf{Partie I --- La transaction.}
\begin{enumerate}
\item Distinguer les \omterm{def:mx:portfolio-trades-and-risk-bids:agency}{transactions de portefeuille en agence} et pour compte propre.
\item Définir un \omterm{def:mx:portfolio-trades-and-risk-bids:agency}{risk bid} et un \omterm{def:mx:portfolio-trades-and-risk-bids:blind}{risk bid à l'aveugle}.
\item Décrire l'univers simulé et le panier.
\item Qu'ont établi Kavajecz et Keim sur les enchères à l'aveugle pour compte propre~?
\end{enumerate}

\medskip\noindent\textbf{Partie II --- Coût et risque.}
\begin{enumerate}[resume]
\item Énoncer la règle de liquidation et donner le coût du panier.
\item Dériver le recouvrement $I(a,b)$ de deux positions à liquidation linéaire.
\item Donner le risque du panier avec et sans la couverture de marché.
\item Que montre le profil de liquidité~?
\end{enumerate}

\medskip\noindent\textbf{Partie III --- Concurrence.}
\begin{enumerate}[resume]
\item Expliquer la malédiction du vainqueur et en dériver la décote.
\item Quel profit les enchérisseurs réalisent-ils sans décote et avec elle~?
\item Comment l'enchérisseur à l'aveugle estime-t-il le coût et le risque, et que constate-t-il~?
\item Donner les offres divulguée et à l'aveugle composante par composante.
\item \emph{Énoncer le résultat nommé}~: l'offre en points de base pour le panier divulgué et pour le panier à l'aveugle, et l'ajustement de malédiction du vainqueur face à trois
concurrents.
\end{enumerate}

\medskip\noindent\textbf{Partie IV --- Transitions et jugement.}
\begin{enumerate}[resume]
\item Chiffrer la transition en bloc, nette et avec croisement interne.
\item Pourquoi le croisement interne est-il moins coûteux que toute exécution~?
\item Quand une offre à l'aveugle vaut-elle sa prime pour le client~?
\item Quelle hypothèse de la tarification vérifieriez-vous en premier~?
\item Comment tarifer l'offre si les concurrents partageaient votre modèle~?
\item Que diriez-vous au fonds sur l'agence par opposition au principal~?
\item En une phrase~: que vend une \omterm{def:mx:portfolio-trades-and-risk-bids:agency}{risk bid}~?
\end{enumerate}
\end{problem}

\section{Questions d'entretien}

\begin{interviewq}[$\star$ \iqroles{trader}]\label{iq:mx:portfolio-trades-and-risk-bids:1}
Un client vous demande un prix de risque sur un panier à la clôture. Que devez-vous savoir, et comment construisez-vous le prix~?
\end{interviewq}

\begin{interviewq}[$\star\star$ \iqroles{chercheur}]\label{iq:mx:portfolio-trades-and-risk-bids:2}
Expliquez la malédiction du vainqueur dans une enchère de \omterm{def:mx:portfolio-trades-and-risk-bids:agency}{risk bids} et la manière dont vous la corrigeriez.
\end{interviewq}

\begin{interviewq}[$\star\star$ \iqroles{risque}]\label{iq:mx:portfolio-trades-and-risk-bids:3}
Comment mesureriez-vous le risque de détenir un panier pendant son débouclage, et que pouvez-vous couvrir~?
\end{interviewq}

\begin{interviewq}[$\star\star$ \iqroles{banque}]\label{iq:mx:portfolio-trades-and-risk-bids:4}
Votre desk continue de perdre de l'argent sur les \omterm{def:mx:portfolio-trades-and-risk-bids:blind}{risk bids à l'aveugle} qu'il remporte. Que vérifiez-vous~?
\end{interviewq}

\begin{interviewq}[$\star\star$ \iqroles{trader}]\label{iq:mx:portfolio-trades-and-risk-bids:5}
Planifiez une transition d'un gérant actions à un autre pour un fonds de pension.
\end{interviewq}

\begin{interviewq}[$\star\star\star$ \iqroles{chercheur}]\label{iq:mx:portfolio-trades-and-risk-bids:6}
Comment estimeriez-vous le coût d'un panier que vous ne pouvez pas voir à partir de son profil de liquidité, et comment sauriez-vous si l'estimation est biaisée~?
\end{interviewq}
